\section{Overall Architecture}
\label{sec:architecture}

\subsection{Layered organization}

\modsim{} is distributed as \code{modsim-robotics} and currently supports Python 3.11
and newer. The source tree is divided into three packages:

\begin{itemize}
  \item \pathcode{src/modsim}: backend-neutral schemas, loading, validation, canonical
        runtime, docking, model views, scenarios, CLI, and mock backend;
  \item \pathcode{src/modsim_backend_mujoco}: optional scene composition, adapter,
        viewer, runtime host process, and weld/contact mechanics; and
  \item \pathcode{src/modsim_studio}: optional Qt/PyVista authoring and Qt/PyQtGraph
        runtime visualization.
\end{itemize}

\Cref{fig:layers} shows the authority and dependency direction. Upward clients depend
on immutable or typed core interfaces; the core does not import optional UI or physics
packages.

\begin{figure}[htbp]
\centering
\begin{tikzpicture}[
  node distance=5mm and 8mm,
  box/.style={draw,rounded corners,minimum height=9mm,align=center,font=\small,
              fill=modsimgray,text width=0.28\textwidth},
  core/.style={box,fill=modsimblue!12,draw=modsimblue},
  opt/.style={box,fill=modsimcyan!12,draw=modsimcyan!80!black},
  data/.style={box,fill=modsimgold!20,draw=modsimgold!70!black},
  arrow/.style={-{Latex[length=2mm]},thick,draw=black!65}
]
  \node[data] (assets) {Mechanical assets\\URDF / MJCF / meshes};
  \node[data,right=of assets] (pack) {Robot Pack\\semantics / mappings / recipes};
  \node[core,below=of $(assets)!0.5!(pack)$,text width=0.60\textwidth] (load)
    {Loader, validator, writer, importer, scene specification};
  \node[core,below=of load,text width=0.60\textwidth] (runtime)
    {Canonical runtime: WorldState, docking, assemblies, events, metrics, commands};
  \node[opt,below left=of runtime] (backends) {Backend adapters\\mock / MuJoCo / future};
  \node[core,below=of runtime] (views) {ModelViewFactory\\immutable derived views};
  \node[opt,below left=of views] (studio) {Robot Pack Studio\\authoring};
  \node[opt,below right=of views] (inspector) {Runtime Inspector\\graph / events / metrics};
  \draw[arrow] (assets) -- (load);
  \draw[arrow] (pack) -- (load);
  \draw[arrow] (load) -- (runtime);
  \draw[arrow] (runtime) -- (backends);
  \draw[arrow] (backends) -- node[below,sloped,font=\scriptsize]{snapshots / outcomes} (runtime);
  \draw[arrow] (runtime) -- (views);
  \draw[arrow] (load) -| (studio);
  \draw[arrow] (views) -- (inspector);
\end{tikzpicture}
\caption{Layered \modsim{} architecture. Arrow direction denotes the primary flow of
data or requests, not Python import direction. Optional physics and UI packages remain
outside the dependency-light core.}
\label{fig:layers}
\end{figure}

\subsection{End-to-end data flow}

An ordinary workflow moves through six stages.

\begin{enumerate}
  \item The user imports a concrete URDF or opens an existing pack. The loader resolves
        split YAML documents and local assets under a pack root.
  \item Validation checks schema and semantic cross-references under an authoring or
        simulation-readiness profile.
  \item A \code{SceneSpec} instantiates module types with stable instance IDs and initial
        poses.
  \item A backend loads the selected mechanical asset for each instance and returns a
        handle registry. The runtime immediately ingests a snapshot.
  \item Each step advances physics, ingests observed state, processes overload and
        docking, appends events, and updates revision counters.
  \item A model-view factory builds an immutable topology projection, and the Runtime
        Inspector publishes it with event deltas, metrics, and scenario status.
\end{enumerate}

The crucial property is that no stage invents a second canonical topology. The backend
reports constraints and frames; the world stores logical connections; the graph is
derived from those connections.

\subsection{Stable identity}

Typed string identifiers are used for pack entities and runtime instances. A connector
instance ID combines a module instance and local connector, for example
\code{module\_3/pan}. A joint instance similarly combines module and semantic joint.
A connection ID is canonical with respect to endpoint ordering, so requesting the same
undirected pair in either order yields the same identity. Assembly IDs derive from
membership, using the smallest module ID in the component, rather than from a process
counter. These rules support deterministic logging, graph keys, replay, and
cross-process transport.

\subsection{Capability-oriented extension}

A backend advertises exactly what it implements. The base contract includes loading,
stepping, snapshots, physical connection creation/removal, shutdown, and capabilities.
Optional protocols add module pose/twist control or joint commands. Core code checks
capabilities before use. For example, a backend without runtime constraints produces a
failed dock rather than a logical-only edge; a backend without effort control rejects a
physical wheel scenario during preflight, before the initial world is mutated.

Model-view extensibility follows the same pattern. A pack recipe names a stable builder
ID, not a Python import path. An application registers builder implementations in a
\code{ModelViewFactory}. Missing builders are valid portable data but unavailable in
that installation, producing an explicit construction error.

\subsection{Package and distribution boundary}

The wheel contains the three Python packages. The source distribution additionally
contains tests, examples, and documentation. This creates one known portability
boundary: the named \smoresep{} demonstration configurations currently live in
\pathcode{examples/scenarios} and are intended for a source checkout. The installed
library exposes the generic scenario machinery, but a wheel without the example tree
cannot execute those named source-checkout demos. A future scenario resource or
explicit \code{--scenario} mechanism should remove this mismatch.

\subsection{Implementation scale}

At the current snapshot, the three source packages contain 72 Python modules and about
20,455 lines of Python. The test directory contains 40 test modules and about 11,073
lines. The committed \smoresep{} pack occupies approximately 16 MiB, dominated by five
binary STL visual meshes. These numbers characterize engineering scope, not runtime
performance; \Cref{sec:analysis} treats performance separately.

